\documentclass[conference]{IEEEtran}
% ---------------------------------------------------------------------------
% Build: pdflatex main && bibtex main && pdflatex main && pdflatex main
% (or upload the whole paper/ folder to Overleaf). Every number in this paper
% is generated by src/analysis.py into latex/numbers.tex and latex/tab_*.tex;
% do not edit numbers by hand -- rerun `python -m src.analysis` instead.
% ---------------------------------------------------------------------------
\IEEEoverridecommandlockouts
\usepackage{cite}
\usepackage{amsmath,amssymb}
\usepackage{graphicx}
\usepackage{booktabs}
\usepackage{multirow}
\usepackage{xcolor}
\usepackage[hidelinks]{hyperref}
\usepackage{url}

\graphicspath{{figures/}}
\input{latex/numbers.tex}

% Double-blind switch: set to \blindtrue for submission, \blindfalse for camera-ready.
\newif\ifblind
\blindtrue

\begin{document}

\title{How Much Does Lesion Leakage Inflate HAM10000 Results?\\
A Reproducible Lesion-Level Benchmark of Lightweight Transfer Learning for Seven-Class Dermoscopic Classification}

\ifblind
\author{\IEEEauthorblockN{Anonymous Authors}\IEEEauthorblockA{Paper under double-blind review}}
\else
\author{\IEEEauthorblockN{Author Name}
\IEEEauthorblockA{\textit{Department} \\
\textit{Institution}\\
City, Country \\
email@example.com}}
\fi

\maketitle

% ===========================================================================
\begin{abstract}
%ABSTRACT%
\end{abstract}

\begin{IEEEkeywords}
skin lesion classification, dermoscopy, HAM10000, data leakage, transfer learning, MobileNetV2, calibration, reproducibility
\end{IEEEkeywords}

% ===========================================================================
\section{Introduction}

Skin cancer is among the most frequently diagnosed cancers worldwide, and early recognition of melanoma in particular strongly influences prognosis. Dermoscopy improves diagnostic accuracy over naked-eye examination, but its interpretation requires training and experience. Convolutional neural networks (CNNs) have been shown to classify dermoscopic and clinical images at a level comparable to dermatologists in reader studies~\cite{esteva2017dermatologist,haenssle2018man,brinker2019deep,tschandl2019comparison}, and public benchmarks such as the ISIC challenges~\cite{codella2018isic2017,codella2019isic2018} have made the task a standard testbed for medical image classification.

The HAM10000 dataset~\cite{tschandl2018ham10000} is the most widely used public benchmark for seven-class pigmented-lesion classification. It contains \nImg{} dermoscopic images but only \nLes{} distinct lesions: many lesions were photographed more than once (different magnifications, angles or time points). If a dataset is partitioned \emph{by image}, images of the same lesion can fall into both the training and the test set, so that the test set partly measures recognition of already-seen lesions rather than generalisation to new ones. This is a textbook instance of data leakage~\cite{kapoor2023leakage,roberts2021pitfalls}, and duplicate or near-duplicate images across ISIC partitions have been documented before~\cite{cassidy2022isic}. Nevertheless, many HAM10000 studies do not report how their partitions were formed, which makes published numbers difficult to compare.

A second gap concerns reporting. Results are often given as a single accuracy value from a single run, although HAM10000 is heavily imbalanced (\,$\approx$67\% melanocytic nevi) so that accuracy is dominated by the majority class, and small classes such as dermatofibroma contribute only a few dozen test images, so that per-class estimates carry wide uncertainty. Reporting guidelines for clinical prediction models~\cite{collins2024tripodai} recommend uncertainty estimates, calibration and clinically meaningful operating points.

This paper provides a small but careful, fully reproducible study that addresses both gaps with a lightweight model that can be trained on a laptop CPU. Our contributions are:
\begin{enumerate}
  \item A \textbf{leakage-free, class-stratified, lesion-grouped} 70/15/15 partition of all \nImg{} HAM10000 images, released as a CSV file together with automated checks that no lesion is shared between partitions.
  \item A \textbf{direct measurement of leakage optimism}: the identical training recipe is evaluated under a lesion-level and a naive image-level split, and the leaky test set is decomposed into images of seen and unseen lesions.
  \item A \textbf{transparent MobileNetV2 baseline} (2.55\,M parameters) with ablations of backbone fine-tuning and class-weighted loss, multiple random seeds, lesion-level bootstrap confidence intervals, paired significance tests, probability calibration, a malignant-versus-benign triage operating point, and Grad-CAM visualisations.
  \item \textbf{Open code} that reproduces every table and figure from raw data on a consumer CPU.
\end{enumerate}

% ===========================================================================
\section{Related Work}

\textit{Deep learning for dermoscopy.} Esteva et al.~\cite{esteva2017dermatologist} trained an Inception-based CNN on about 129{,}000 clinical images and reported dermatologist-level performance on biopsy-proven test sets. Haenssle et al.~\cite{haenssle2018man} and Brinker et al.~\cite{brinker2019deep} compared CNNs with 58 and 157 dermatologists, respectively, for dermoscopic melanoma recognition, and Tschandl et al.~\cite{tschandl2019comparison} showed that the best ISIC~2018 algorithms outperformed human readers on average in a large web-based study, while also noting reduced performance on out-of-distribution images.

\textit{Benchmarks.} HAM10000~\cite{tschandl2018ham10000} was released as the training set of ISIC~2018 Task~3~\cite{codella2019isic2018}, whose official ranking metric is balanced multi-class accuracy (mean per-class recall) precisely because of class imbalance. Top challenge entries typically combined large ensembles of high-capacity networks~\cite{he2016resnet,tan2019efficientnet}, higher input resolutions and, in some cases, external data and metadata~\cite{gessert2020ensembles}. Such systems are not directly comparable to single lightweight models, and they are not the focus of this work. Other public dermoscopic datasets include BCN20000~\cite{combalia2019bcn20000} and PH$^2$~\cite{mendonca2013ph2}.

\textit{Leakage, shortcuts and bias.} Kapoor and Narayanan~\cite{kapoor2023leakage} catalogue leakage as a pervasive cause of over-optimistic results in ML-based science, and Roberts et al.~\cite{roberts2021pitfalls} found similar methodological pitfalls, including patient overlap between training and test data, across hundreds of medical-imaging models. Cassidy et al.~\cite{cassidy2022isic} analysed ISIC datasets and reported duplicate images and the need for careful partitioning. Winkler et al.~\cite{winkler2019markings} demonstrated that a CNN's melanoma scores were influenced by surgical skin markings, a spurious shortcut. Adamson and Smith~\cite{adamson2018disparities}, Groh et al.~\cite{groh2021fitzpatrick} and Daneshjou et al.~\cite{daneshjou2022disparities} highlight that public dermatology datasets under-represent darker skin tones and that model performance can drop on diverse populations.

\textit{Imbalance and calibration.} Class re-weighting~\cite{cui2019classbalanced} and focal loss~\cite{lin2017focal} are common remedies for long-tailed label distributions. Guo et al.~\cite{guo2017calibration} showed that modern CNNs are often mis-calibrated and that a single temperature parameter fitted on held-out data is an effective correction.

% ===========================================================================
\section{Materials and Methods}

\subsection{Dataset}
We use all \nImg{} images of HAM10000~\cite{tschandl2018ham10000} (CC BY-NC 4.0) with the seven diagnostic labels of the official metadata: actinic keratosis/intra-epithelial carcinoma (\texttt{akiec}), basal cell carcinoma (\texttt{bcc}), benign keratosis-like lesions (\texttt{bkl}), dermatofibroma (\texttt{df}), melanoma (\texttt{mel}), melanocytic nevus (\texttt{nv}) and vascular lesions (\texttt{vasc}). The ground truth was established by histopathology for 53\% of the images and otherwise by follow-up examination, expert consensus or confocal microscopy~\cite{tschandl2018ham10000}. The images (600$\times$450 pixels) were downloaded from the public ISIC archive; no image was excluded. For the binary triage analysis, \texttt{mel}, \texttt{bcc} and \texttt{akiec} are grouped as malignant or pre-malignant and the remaining classes as benign.

\subsection{Lesion-level partitioning}
Each image carries a \texttt{lesion\_id}. We assign whole lesions to 20 folds with a greedy class-stratified group assignment (in the spirit of stratified group $k$-fold): lesions are processed from the most to the least class-specific and each is placed in the fold whose class distribution deviates least from the global one. Three folds form the test set, three the validation set and fourteen the training set (70/15/15). Programmatic assertions verify that the three partitions share no lesion and that every class occurs in every partition. The resulting partition (Table~\ref{tab:dataset}) contains \nTrainImg{} / \nValImg{} / \nTestImg{} images from \nTrainLes{} / \nValLes{} / \nTestLes{} lesions and is released with the code.

\begin{table}[t]
\centering
\caption{HAM10000 class distribution in the lesion-level partition (images). No lesion appears in more than one partition.}
\label{tab:dataset}
\resizebox{\columnwidth}{!}{\input{latex/tab_dataset.tex}}
\end{table}

For the leakage experiment only, we additionally draw a naive class-stratified 70/15/15 split \emph{by image} with the same proportions (seed 42), as is common in the literature. Under this split, \leakSeenN{} of the \leakImgN{} test images (\leakSeenPct\%) belong to a lesion that also has at least one image in the training set.

\subsection{Model}
We use MobileNetV2~\cite{sandler2018mobilenetv2} pre-trained on ImageNet~\cite{deng2009imagenet} (torchvision weights \texttt{IMAGENET1K\_V1}). The original classifier is replaced by a head of dropout ($p{=}0.4$), a 256-unit linear layer with ReLU, dropout ($p{=}0.3$) and a 7-way linear output. The model has 2.55\,M parameters, of which 0.33\,M are in the head. We compare two regimes: (i) \emph{frozen}, where only the head is trained and the backbone's batch-normalisation statistics are kept fixed; and (ii) \emph{fine-tuned}, where all layers are trained with a ten-times smaller learning rate for the backbone than for the head.

\subsection{Training}
Images are resized to 224$\times$224 and normalised with ImageNet statistics. Because dermoscopic images have no canonical orientation, training-time augmentation applies random horizontal and vertical flips, a rotation by an arbitrary angle, isotropic scaling by a factor in $[0.83, 1.18]$, a small translation (up to $\approx$4\% of the image width) with reflection padding, and brightness and contrast jitter of $\pm$15\%. We use AdamW~\cite{loshchilov2019adamw} (weight decay $10^{-4}$) with learning rates of $10^{-3}$ (head) and $10^{-4}$ (backbone), batch size 32, and 10 epochs with a one-epoch linear warm-up followed by cosine decay~\cite{loshchilov2017sgdr}. Training uses no oversampling.

\textit{Class-weighted loss.} In the weighted configurations the cross-entropy loss weights each class by the square root of its inverse training frequency, $w_c \propto (N / (K n_c))^{1/2}$, normalised to mean one (resulting weights: \texttt{akiec} 1.10, \texttt{bcc} 0.88, \texttt{bkl} 0.60, \texttt{df} 1.88, \texttt{mel} 0.60, \texttt{nv} 0.24, \texttt{vasc} 1.69). The square root tempers the full inverse-frequency weights, which over-emphasise the smallest classes~\cite{cui2019classbalanced}.

\textit{Model selection.} After every epoch we compute balanced accuracy on the validation set and keep the best checkpoint. The test set is used exactly once per trained model, after training has finished.

\textit{Configurations.} We train four configurations: frozen or fine-tuned backbone, each with plain or class-weighted cross-entropy. The primary configuration (fine-tuned, class-weighted) is trained with \nSeeds{} random seeds; ablations with one seed. For the leakage experiment the primary recipe is trained once more on the image-level split.

\subsection{Evaluation}
We report accuracy; balanced accuracy (mean per-class recall, the ISIC~2018 ranking metric~\cite{codella2019isic2018}); macro- and support-weighted F1; one-vs-rest ROC AUC per class and its macro average; melanoma sensitivity; and, for the binary malignant-versus-benign grouping, sensitivity and specificity of the arg-max prediction and the AUC of the summed malignant probability.

\textit{Uncertainty.} Because images of one lesion are correlated, 95\% confidence intervals are obtained with a \emph{lesion-level} percentile bootstrap~\cite{efron1993bootstrap}: 1{,}000 times, test lesions are resampled with replacement and all of their images are included. Differences between configurations are assessed with a paired lesion-level bootstrap on the same test images; we report the difference, its 95\% CI and a two-sided bootstrap $p$-value. Across seeds we report mean $\pm$ standard deviation, and the detailed per-class results of the primary configuration use the average of its seeds' predicted probabilities (a seed ensemble).

\textit{Calibration and triage operating point.} We measure the expected calibration error (ECE, 15 equal-width bins) and fit a single softmax temperature on the validation set~\cite{guo2017calibration}. To emulate a triage setting in which missing a malignancy is costlier than a false alarm, we choose, on the validation set, the threshold on the summed malignant probability that reaches 90\% (and 95\%) malignant sensitivity, and then report sensitivity, specificity and the fraction of flagged cases on the test set.

\textit{Explainability.} We visualise class-discriminative regions with Grad-CAM~\cite{selvaraju2017gradcam} computed on the last convolutional feature map for randomly chosen test images of every class.

\subsection{Implementation}
All experiments use PyTorch~\cite{paszke2019pytorch} on a laptop CPU (Intel Core i7-14650HX, 16 threads, no GPU). Metrics are implemented in NumPy and unit-tested against hand-computed values. The code, partition file and scripts that regenerate every number, table and figure are publicly available%CODELINK%.

% ===========================================================================
\section{Results}
%RESULTS%

% ===========================================================================
\section{Discussion}
%DISCUSSION%

% ===========================================================================
\section{Conclusion}
%CONCLUSION%

% ===========================================================================
\section*{Ethics Statement}
This study uses only the publicly available, de-identified HAM10000 dataset~\cite{tschandl2018ham10000}, released under a CC BY-NC 4.0 licence; its collection was approved by the ethics committees of the contributing institutions~\cite{tschandl2018ham10000}. No new patient data were collected, and no additional ethical approval was required. The models described here are research prototypes and are not intended for clinical use.

\section*{Data and Code Availability}
HAM10000 is available from the ISIC archive and the Harvard Dataverse (doi:10.7910/DVN/DBW86T). The lesion-level partition, training and evaluation code, and all raw prediction files are available%CODELINK2%.

\bibliographystyle{IEEEtran}
\bibliography{references}

\end{document}
